\documentclass{beamer}
\usepackage[italian]{babel}
\usepackage[utf8]{inputenc}
\usepackage{graphicx}
\usepackage{subfig}
\usepackage{hyperref}
\usepackage{tikz}
\usetikzlibrary{shapes, calc, positioning}
\usepackage{tabularx}
\usepackage{listings}
\usepackage{color}

\definecolor{mygreen}{rgb}{0,0.6,0}
\definecolor{mygray}{rgb}{0.5,0.5,0.5}
\definecolor{mymauve}{rgb}{0.58,0,0.82}

\lstset{ 
  backgroundcolor=\color{black!20!white},   % choose the background color
  basicstyle=\ttfamily,        % the size of the fonts that are used for the code
  captionpos=b,                    % sets the caption-position to bottom
  commentstyle=\color{mygreen},    % comment style
  frame=single,	                   % adds a frame around the code
  %keepspaces=true,                 % useful for keeping indentation 
  language=bash,
  rulecolor=\color{black},         % if not set, the frame-color may be changed on line-breaks within not-black text (e.g. comments (green here))
  showspaces=false,                % show spaces everywhere adding particular underscores; it overrides 'showstringspaces'
  showtabs=false,                  % show tabs within strings adding particular underscores
  tabsize=2,	                   % sets default tabsize to 2 spaces
    deletekeywords={cd},            % if you want to delete keywords from the given language
}

\title{Terminale avanzato}
\author{Samuele Musiani, Alice Benatti}
\institute{Università di Bologna, corso di Laurea in Informatica}
\date{12 ottobre 2023}
\logo{\includegraphics[width=0.05\textwidth]{assets/by-nc-sa-4-0.png}}

\AtBeginSection[]{
  \begin{frame}
    % \setcounter{tocdepth}{10}
  \vfill
  \centering
  \begin{beamercolorbox}[sep=8pt,center,shadow=true,rounded=true]{title}
    \usebeamerfont{title}\insertsectionhead\par%
  \end{beamercolorbox}
  % \tableofcontents[currentsubsection]
  \vfill
\end{frame}
}

\begin{document}

\begin{frame} 
  \titlepage
\end{frame}

\section{Un mare di shell}

\subsection{Cos'è una shell}
\begin{frame}[fragile]{Cos'è una shell}
  Una \texttt{shell} è un programma che permette di parlare con il sistema 
  operativo attraverso dei comandi da tastiera.
  \pause

  \begin{figure}
        \begin{lstlisting}
[samu@leibniz ~]$
    \end{lstlisting}
    \caption{Esempio di un \textit{shell prompt}}
  \end{figure}
\end{frame}

\subsection{Cos'è una shell}
\begin{frame}[fragile]{Cos'è una shell}
  Esistono varie shell:
  \begin{itemize}
    \item sh
    \item bash
    \item zsh
    \item fish
  \end{itemize}
  Tutte servono per interagire con il sistema operativo, però sono tutte 
  leggermente diverse sotto alcuni aspetti.\medskip

  Di solito la predefinita nei sistemi Linux è bash.
\end{frame}

\section{I primi passi in un terminale}
\begin{frame}{Aprire un terminale}
  Per utilizzare una \texttt{shell} è necessario disporre di un emulatore di 
  terminale.\bigskip

  Un emulatore di terminale è un'applicazione che permette di interagire con
  la \texttt{shell}.\medskip

  Esisto vari tipi di emulatori che offrono configurazioni e opzioni diverse:
  \begin{itemize}
    \item Alacritty
    \item Foot
    \item Kitty
    \item Konsole
    \item ...
  \end{itemize}
\end{frame}

\section{Pipe e ridirezione: il potere della shell}
\begin{frame}{Il potere della shell}
  Di solito si utilizzano comandi singoli. Nonostante molti di essi siano utili 
  anche da soli, sono sicuramente più utili usati in combinazione con altri 
  comandi.\bigskip

  Come si fa però a combinare più comandi? \pause
  Ricordiamoci che:
  \begin{itemize}
    \item <2-> Un comando restituisce sempre qualcosa sullo standard output
    \item <3-> Lo standard output è considerato come un file (in memoria) dal 
      sistema 
    \item <4-> La maggior parte dei comandi visti fino ad adesso hanno la 
      possibilità di prendere in input lo standard output invece che un file 
      classico
  \end{itemize}
\end{frame}

\begin{frame}{Il potere della shell}
  \begin{figure}
    \begin{tikzpicture}
      \draw (-1.5,3) rectangle ++(3, 1);
      \draw (-1.5,-3) rectangle ++(3, 1);
      \draw [->] (0,3) -- (0, 1);
      \draw [->] (0,-0.9) -- (0, -2);
      \node [cylinder, shape border rotate=90, draw,minimum height=2cm,minimum width=1.5cm] {};

      \node {buffer};
      \node at (0,3.5) {command 1};
      \node at (0,-2.5) {command 2};
      \node at (1, 2) {output};
      \node at (1,-1.5) {input};
    \end{tikzpicture}
  \end{figure}
\end{frame}

\begin{frame}{Pipe}
  Per prendere l'output di un comando e riderizionarlo in input verso un altro
  comando si usa la \textbf{pipe} \texttt{|}\bigskip

  Per esempio se vogliamo vedere tutti i file presenti in \texttt{/bin} il 
  nostro terminale si riempie di scritte.\bigskip

  Possiamo visualizzare il lungo output con il comando \texttt{less}: 
  \texttt{ls /bin | less}
\end{frame}

\begin{frame}{Pipe - grep}
  Come abbiamo visto la lista di file presenti in \texttt{/bin} è molto lunga.
  Se volessimo trovarne uno specifico?\pause \bigskip

  Nonostante esiste un comando apposito per cercare file, possiamo fare:
  \texttt{ls /bin | grep "firefox"} dove al posto di \textit{firefox} può 
  andarci una qualsiasi stringa.
\end{frame}

\begin{frame}{Pipe - esempi}
  Di seguito una lista di esempi di utilizzo della \texttt{pipe}:
  \begin{itemize}
    \item \texttt{cat file1 file2 | grep "word"} cerca una stringa in più file
    \item \texttt{ls /bin | wc -l} conta quanti programmi sono presenti in 
      \texttt{/bin}
    \item \texttt{ls /bin | grep "zip" | wc -l} conta quanti programmi hanno
      la stringa "\textit{zip}" al loro interno nella cartella \texttt{/bin}
    \item \texttt{grep "castoro" animali | wc -l} conta le occorrenze di 
      \texttt{castoro} trovate nel file \texttt{animali}
    \item \texttt{grep "the" book | less} mostra le occorrenze di \texttt{the} 
      trovate in \texttt{book} attraverso il lettore \texttt{less}
  \end{itemize}
\end{frame}

\begin{frame}{Ridirezione su file}
  Come abbiamo visto è possibile mandare l'output di un comando nell'input di
  un altro comando.\bigskip

  Se volessimo salvare l'output di un comando su un file?\pause

  Esiste l'\textbf{operatore di ridirezione} \texttt{>}\bigskip

  Al posto di indirizzare l'output in un comando, scrive direttamente su un
  file.\bigskip

  Sintassi: \texttt{comando > file}\bigskip

  ATTENZIONE: Alla \textit{shell} non interessa se il file esiste già, quindi se 
  esiste lo SOVRASCRIVE COMPLETAMENTE.
\end{frame}

\begin{frame}{Ridirezione su file non distruttiva}
  Esiste anche un operatore per indirizzare su file l'output di un comando senza 
  sovrascrivere il contenuto del file, ma "appendendo" alla fine del file il
  contenuto scritto.\bigskip

  Sintassi: \texttt{comando >> file}\bigskip

  Si usa nello stesso modo dell'operatore classico\bigskip
\end{frame}

\begin{frame}[fragile]{Ridirezione dello standard error}
  Di predefinito l'operatore di redirezione, come la pipe, reindirizza solo
  lo standard output. Questo significa che se un comando restituisce un errore
  questo non verrà riportato nel file indicato:

  \begin{figure}
    \begin{lstlisting}
[samu@leibniz ~]$ ls /asdkjc > /tmp/castoro
ls: cannot access '/asdkjc': No such file or 
directory
[samu@leibniz ~]$
    \end{lstlisting}
    \caption{Esempio di stderr non ridirezionato}
  \end{figure}

  Per ridirezionare anche lo standard error dentro il file è necessario
  specificarlo alla shell.
\end{frame}

\begin{frame}[fragile]{Ridirezione dello standard error}
  Lo standard error in linux ha come \textit{file descriptor} il \textbf{2}. 
  Non importa sapere cosa sia un file descriptor, basta ricordarsi che
  \textbf{stdout = 1} e \textbf{stderr = 2}. \medskip

  Lo standard error si può redirezionare solo verso un altro file descriptor.
  Quindi per ridirezionare lo stderr nello stdout si usa la seguente sintassi:
  \textbf{2>\&1} \medskip

  Redirezionando lo standard error verso lo standard output è possibile scrivere
  entrambi su file:
  \begin{figure}
    \begin{lstlisting}
[samu@leibniz ~]$ ls /asdkjc 2>&1 > /tmp/test
[samu@leibniz ~]$ cat /tmp/test
ls: cannot access '/asdkjc': No such file or 
directory
    \end{lstlisting}
    \caption{Esempio di stderr ridirezionato}
  \end{figure}
\end{frame}

\begin{frame}[fragile]{Ridirezione dello standard error}
  Da notare che questo funziona anche per la pipe:
  \begin{figure}
    \begin{lstlisting}
[samu@leibniz ~]$ ls /asdkjc | head -c 10
ls: cannot access '/asdkjc': No such file or 
directory
[samu@leibniz ~]$ ls /asdkjc 2>&1 | head -c 10
ls: cannot
    \end{lstlisting}
    \caption{Esempio di stderr ridirezionato}
  \end{figure}
\end{frame}

\begin{frame}{Filtri}
  Quando si combinano assieme più comandi spesso i seguenti tornano utili:
  \begin{itemize}
    \item<1-> sort: Restituisce in stdout il testo preso in stdin
    \item<2-> uniq: Se ci sono più linee consecutive uguali ne lascia solo una
    \item<3-> grep: Cerca il testo passato come argomento
    \item<4-> wc: Conta linee, parole e carateri passati in input
    \item<5-> head: Manda in stdout solo la prima parte del testo
    \item<6-> tail: Manda in stdout solo l'ultima parte del testo
  \end{itemize}
\end{frame}

\begin{frame}[fragile]{/dev/null}
  Come si fa a ridirezionare soltanto lo stderr su file e "buttare via" lo
  stdout? \medskip
  \pause

  Esiste in linux un "file" che qualsiasi cosa gli passiate lo butta via. Questo
  "file" è \textbf{/dev/null}.\medskip

  \begin{figure}
    \begin{lstlisting}
[samu@leibniz ~]$ ls /asdkjc / 2> /tmp/pippo 
1> /dev/null
[samu@leibniz ~]$ cat /tmp/pippo
ls: cannot access '/asdkjc': No such file or 
directory
    \end{lstlisting}
    \caption{Esempio di stderr ridirezionato}
  \end{figure}

Ricordiamoci che stdout = 1 e stderr = 2.
\end{frame}

\section{\Large echo, \small echo, \tiny echo, echo...}
